% GENERATED FILE. Edit canonical lessons, then run npm run generate:books.
% canonical-source-hash: fnv1a64:52abe5b87c8a38bc
% canonical-lessons: BN-C492-karigor, BN-C492-sekra, BN-C492-jhalai-kora, BN-C492-kasari, BN-C492-gatha

\chapter{Skilled Hands}
\label{ch:bn-skilled-hands}
\hlchaptermodality{492}

\begin{chapteropening}
\textbf{By the end of this chapter:} \emph{I can name some skilled trades and say what they make and mend.}

Complete the last lesson of chapter 492: I can name some skilled trades and say what they make and mend.
\end{chapteropening}

\section[\texorpdfstring{\bn{কারিগর} (kārigôr)}{কারিগর (kārigôr)}]{\bn{কারিগর} (kārigôr) — a craftsman; a skilled worker}

\label{lesson:BN-C492-karigor}



\begin{quote}
Before the new one: say the Bengali for a female singer, then the Bengali for a female dancer.
\end{quote}

\subsection*{You'll want to know: \bn{কারিগর}}
\textbf{\bn{কারিগর}} — \emph{kārigôr} — \textquotedblleft{}a craftsman; a skilled worker\textquotedblright{}.

\begin{quote}
\bn{আমাদের} \bn{গ্রামের} \bn{কারিগর} \bn{এই} \bn{চেয়ারটা} \bn{বানাল।}
\end{quote}

\begin{quote}
\emph{Āmāder grāmer kārigôr ei cheyārṭā bānālô.}
\end{quote}

\begin{quote}
A craftsman from our village made this chair.
\end{quote}

\begin{cousinweb}
\textbf{\bn{কারিগর}} (\emph{kārigôr}) is usually traced to Persian \emph{kārīgar}, \textquotedblleft{}a craftsman\textquotedblright{}, built on \emph{kār}, \textquotedblleft{}work\textquotedblright{}.
\end{cousinweb}

\begin{grammarlens}[title={What you've built}]
One more word for describing what you lived and felt.
\end{grammarlens}

\subsection*{Your turn}
\begin{itemize}
\raggedright
  \item \emph{Say it:} \emph{kārigôr}
  \item \emph{Say it:} \emph{kārigôr}, once more
  \item \emph{Say it:} \emph{Āmāder grāmer kārigôr ei cheyārṭā bānālô.} — A craftsman from our village made this chair.
  \item \emph{Recall:} say the Bengali for a female singer, then the Bengali for a female dancer, then say \emph{kārigôr} again
\end{itemize}

\begin{tcolorbox}[breakable,colback=teal!4,colframe=teal!35!black,title={Before you move on}]
What does \bn{কারিগর} mean? (\textquotedblleft{}A craftsman; a skilled worker\textquotedblright{}.) Say it once more.
\end{tcolorbox}
\section[\texorpdfstring{\bn{স্যাকরা} (sækrā)}{স্যাকরা (sækrā)}]{\bn{স্যাকরা} (sækrā) — a goldsmith}

\label{lesson:BN-C492-sekra}



\begin{quote}
Before the new one: say the Bengali for a female dancer, then the Bengali for a craftsman.
\end{quote}

\subsection*{You'll want to know: \bn{স্যাকরা}}
\textbf{\bn{স্যাকরা}} — \emph{sækrā} — \textquotedblleft{}a goldsmith\textquotedblright{}.

\begin{quote}
\bn{স্যাকরা} \bn{দিদিমার} \bn{হারটা} \bn{সারাল।}
\end{quote}

\begin{quote}
\emph{Sækrā didimār hārṭā sārālô.}
\end{quote}

\begin{quote}
The goldsmith mended my grandmother's necklace.
\end{quote}

\begin{cousinweb}
\textbf{\bn{স্যাকরা}} (\emph{sækrā}) is usually traced to Sanskrit \emph{svarṇakāra}, \textquotedblleft{}a gold-worker\textquotedblright{} (Monier-Williams), worn down in everyday speech.
\end{cousinweb}

\begin{grammarlens}[title={What you've built}]
One more word for describing what you lived and felt.
\end{grammarlens}

\subsection*{Your turn}
\begin{itemize}
\raggedright
  \item \emph{Say it:} \emph{sækrā}
  \item \emph{Say it:} \emph{sækrā}, once more
  \item \emph{Say it:} \emph{Sækrā didimār hārṭā sārālô.} — The goldsmith mended my grandmother's necklace.
  \item \emph{Recall:} say the Bengali for a female dancer, then the Bengali for a craftsman, then say \emph{sækrā} again
\end{itemize}

\begin{tcolorbox}[breakable,colback=teal!4,colframe=teal!35!black,title={Before you move on}]
What does \bn{স্যাকরা} mean? (\textquotedblleft{}A goldsmith\textquotedblright{}.) Say it once more.
\end{tcolorbox}
\section[\texorpdfstring{\bn{ঝালাই} \bn{করা} (jhālāi kôrā)}{ঝালাই করা (jhālāi kôrā)}]{\bn{ঝালাই} \bn{করা} (jhālāi kôrā) — to solder; to weld}

\label{lesson:BN-C492-jhalai-kora}



\begin{quote}
Before the new one: say the Bengali for a craftsman, then the Bengali for a goldsmith.
\end{quote}

\subsection*{You'll want to know: \bn{ঝালাই} \bn{করা}}
\textbf{\bn{ঝালাই} \bn{করা}} — \emph{jhālāi kôrā} — \textquotedblleft{}to solder; to weld\textquotedblright{}.

\begin{quote}
\bn{মিস্ত্রি} \bn{ভাঙা} \bn{পাইপটা} \bn{ঝালাই} \bn{করল।}
\end{quote}

\begin{quote}
\emph{Mistri bhāngā pāipṭā jhālāi kôrlô.}
\end{quote}

\begin{quote}
The mechanic soldered the broken pipe.
\end{quote}

Past (he/she …): \bn{তিনি} \bn{সাইকেলটা} \bn{ঝালাই} \bn{করলেন}, \emph{tini sāikelṭā jhālāi kôrlen}.

\begin{cousinweb}
\textbf{\bn{ঝালাই} \bn{করা}} (\emph{jhālāi kôrā}) joins \bn{ঝালাই}, \textquotedblleft{}soldering\textquotedblright{}, built on a verb CDIAL lists as Bengali \emph{jhālā}, \textquotedblleft{}to solder\textquotedblright{}, from a reconstructed \emph{jhālayati} (CDIAL 5382), to \bn{করা}, \textquotedblleft{}to do\textquotedblright{}.
\end{cousinweb}

\begin{grammarlens}[title={What you've built}]
One more word for describing what you lived and felt.
\end{grammarlens}

\subsection*{Your turn}
\begin{itemize}
\raggedright
  \item \emph{Say it:} \emph{jhālāi kôrā}
  \item \emph{Say it:} \emph{jhālāi kôrā}, once more
  \item \emph{Say it:} \emph{Mistri bhāngā pāipṭā jhālāi kôrlô.} — The mechanic soldered the broken pipe.
  \item \emph{Recall:} say the Bengali for a craftsman, then the Bengali for a goldsmith, then say \emph{jhālāi kôrā} again
\end{itemize}

\begin{tcolorbox}[breakable,colback=teal!4,colframe=teal!35!black,title={Before you move on}]
What does \bn{ঝালাই} \bn{করা} mean? (\textquotedblleft{}To solder; to weld\textquotedblright{}.) Say it once more.
\end{tcolorbox}
\section[\texorpdfstring{\bn{কাঁসারি} (kā̃sāri)}{কাঁসারি (kā̃sāri)}]{\bn{কাঁসারি} (kā̃sāri) — a bell-metal smith; a brazier}

\label{lesson:BN-C492-kasari}



\begin{quote}
Before the new one: say the Bengali for a goldsmith, then the Bengali for to solder.
\end{quote}

\subsection*{You'll want to know: \bn{কাঁসারি}}
\textbf{\bn{কাঁসারি}} — \emph{kā̃sāri} — \textquotedblleft{}a bell-metal smith; a brazier\textquotedblright{}.

\begin{quote}
\bn{কাঁসারি} \bn{একটা} \bn{নতুন} \bn{থালা} \bn{বানাল।}
\end{quote}

\begin{quote}
\emph{Kā̃sāri ekṭā nôtun thālā bānālô.}
\end{quote}

\begin{quote}
The bell-metal smith made a new plate.
\end{quote}

\begin{cousinweb}
\textbf{\bn{কাঁসারি}} (\emph{kā̃sāri}) is a tadbhava: CDIAL lists Bengali \emph{kā̃sāri}, \textquotedblleft{}a brazier, a coppersmith\textquotedblright{}, under Sanskrit \emph{kāṁsyakāra}, \textquotedblleft{}a worker in bell-metal\textquotedblright{} (CDIAL 2989).
\end{cousinweb}

\begin{grammarlens}[title={What you've built}]
One more word for describing what you lived and felt.
\end{grammarlens}

\subsection*{Your turn}
\begin{itemize}
\raggedright
  \item \emph{Say it:} \emph{kā̃sāri}
  \item \emph{Say it:} \emph{kā̃sāri}, once more
  \item \emph{Say it:} \emph{Kā̃sāri ekṭā nôtun thālā bānālô.} — The bell-metal smith made a new plate.
  \item \emph{Recall:} say the Bengali for a goldsmith, then the Bengali for to solder, then say \emph{kā̃sāri} again
\end{itemize}

\begin{tcolorbox}[breakable,colback=teal!4,colframe=teal!35!black,title={Before you move on}]
What does \bn{কাঁসারি} mean? (\textquotedblleft{}A bell-metal smith; a brazier\textquotedblright{}.) Say it once more.
\end{tcolorbox}
\section[\texorpdfstring{\bn{গাঁথা} (gā̃thā)}{গাঁথা (gā̃thā)}]{\bn{গাঁথা} (gā̃thā) — to string together; to lay (bricks)}

\label{lesson:BN-C492-gatha}



\begin{quote}
Before the new one: say the Bengali for to solder, then the Bengali for a bell-metal smith.
\end{quote}

\subsection*{You'll want to know: \bn{গাঁথা}}
\textbf{\bn{গাঁথা}} — \emph{gā̃thā} — \textquotedblleft{}to string together; to lay (bricks)\textquotedblright{}.

\begin{quote}
\bn{মজুর} \bn{ইট} \bn{গেঁথে} \bn{দেওয়াল} \bn{বানাল।}
\end{quote}

\begin{quote}
\emph{Mojur iṭ gẽthe deoyāl bānālô.}
\end{quote}

\begin{quote}
The labourer laid the bricks and built a wall.
\end{quote}

Past (he/she …): \bn{তিনি} \bn{ফুল} \bn{গাঁথলেন}, \emph{tini phul gā̃thlen}.

The -\bn{এ} form raises the vowel: \bn{গেঁথে}, having laid.

\begin{cousinweb}
\textbf{\bn{গাঁথা}} (\emph{gā̃thā}) is a tadbhava from Sanskrit \emph{granthayati}, \textquotedblleft{}ties, strings together\textquotedblright{}; CDIAL gives Bengali \emph{gā̃thā} for threading and for laying bricks (CDIAL 4353).
\end{cousinweb}

\begin{grammarlens}[title={What you've built}]
One more word for describing what you lived and felt.
\end{grammarlens}

\subsection*{Your turn}
\begin{itemize}
\raggedright
  \item \emph{Say it:} \emph{gā̃thā}
  \item \emph{Say it:} \emph{gā̃thā}, once more
  \item \emph{Say it:} \emph{Mojur iṭ gẽthe deoyāl bānālô.} — The labourer laid the bricks and built a wall.
  \item \emph{Recall:} say the Bengali for to solder, then the Bengali for a bell-metal smith, then say \emph{gā̃thā} again
\end{itemize}

\begin{tcolorbox}[breakable,colback=teal!4,colframe=teal!35!black,title={Before you move on}]
What does \bn{গাঁথা} mean? (\textquotedblleft{}To string together; to lay (bricks)\textquotedblright{}.) Say it once more.
\end{tcolorbox}
